\documentclass[12pt]{article}
\usepackage[T1]{fontenc}
\usepackage[utf8]{inputenc}
\usepackage{lmodern}
\usepackage{textcomp}
\usepackage{enumitem}
\usepackage{geometry}
\usepackage{graphicx}
\usepackage{amsmath, amssymb}
\geometry{margin=1in}
\begin{document}
\begin{center}
Sociology - Undergraduate Level Assessment 6 \
Total Marks: 40 \
Answer all questions.
\end{center}

\section*{Multiple Choice (10 marks)}
\begin{enumerate}[label=\arabic*.]
\item The 'new racism' refers to:
\begin{enumerate}[label=(\alph*)]
    \item a subtler form of prejudice, masked by nationalist pride
    \item a post-modern deconstruction of racist ideas to reveal their lack of depth
    \item racist practices found in newly emerging areas of social life, such as cyberspace
    \item an anti-fascist movement which challenges nationalist politics
\end{enumerate}
\item The 'demographic transition' is a social trend that involves:
\begin{enumerate}[label=(\alph*)]
    \item a reduction in population size, caused by a higher rate of emigration than immigration
    \item a change in the principal causes of death and disease since industrialization
    \item increased birth and death rates, resulting in a relatively young population
    \item a decline in the birth rate, greater life expectancy, and an ageing population
\end{enumerate}
\item Traditional working class identity was based around:
\begin{enumerate}[label=(\alph*)]
    \item shared working conditions in the manufacturing industry
    \item the class consciousness of members of the proletariat
    \item local communities, extended kinship networks and shared leisure pursuits
    \item collective aspirations to move into the middle class
\end{enumerate}
\item Lombroso claimed that:
\begin{enumerate}[label=(\alph*)]
    \item criminals were socialized into an 'underworld' of crime
    \item no act is intrinsically deviant
    \item biological failings drove some people into crime
    \item women were less likely to be arrested than men
\end{enumerate}
\item Marriage appears to be in decline because:
\begin{enumerate}[label=(\alph*)]
    \item the proportion of people living alone has fallen to 29\%
    \item many people are cohabiting in long term relationships
    \item the upward curve of remarriages compensates for the drop in first marriages
    \item all of the above
\end{enumerate}
\end{enumerate}

\section*{True / False (10 marks)}
\textit{Answer True or False}
\begin{enumerate}[label=\arabic*.]
\setcounter{enumi}{5}
\item True or False: The ruling elite is composed exclusively of individuals from the upper class with no representation from other class backgrounds.
\item True or False: Cultural restructuring involves regenerating cities in economic decline, transforming industrial landscapes, and leveraging media advertising.
\item True or False: Marx predicted that class polarization results in a growing gap between rich and poor, fostering class consciousness among the working class.
\item True or False: The ecological approach to urban sociology focuses on the forms of wildlife and natural habitats found in urban areas.
\item True or False: Ethnographic research produces qualitative data because it compares findings from a number of different cases.
\end{enumerate}

\section*{Long Form Response (20 marks)}
\textit{Answer each question with a clear and complete explanation.}
\begin{enumerate}[label=\arabic*.]
\setcounter{enumi}{10}
\item Discuss the key processes of decision-making in the USA as identified by Domhoff, and critically analyze why the exploitation process is not included in his framework.
\item Discuss the concept of social stratification in sociology, analyzing how it affects individuals' movement within societal layers and its implications for social inequality.
\end{enumerate}

\vfill
\noindent\textit{End of Paper}
\end{document}